\chapter{Équations différentielles stochastiques}\label{ch:qm:stochastic-differential-equations}

À 02:00, le calcul de risque de nuit s'arrête. Un processus de variance en racine carrée, avancé
pas à pas, jour après jour, avec le schéma d'Euler simple, a produit une variance négative~; le code en a
pris la racine carrée, et un NaN s'est propagé d'une trajectoire à la valeur du livre et à toutes les limites qui
en dépendent. Ce n'était pas une trajectoire malchanceuse~: avec les paramètres calibrés du desk, trois
trajectoires sur quatre deviennent négatives en moins d'un an, et diviser quatre fois le pas de temps par deux ne change rien.
Le modèle n'était pas faux, c'était le schéma~; le nombre qui le dit, le ratio de Feller, est une
fonction tenant en une ligne des paramètres. Ce chapitre présente ce qu'est une \omterm{def:qm:stochastic-differential-equations:sde}{équation différentielle
stochastique} et quand elle admet une solution, les trois diffusions que la série utilise le plus (géométrique,
à \omterm{def:qm:stochastic-differential-equations:ou}{retour à la moyenne} et en racine carrée), le générateur qui les relie aux équations aux dérivées
partielles, et la formule de Feynman--Kac qui transforme une espérance en EDP et inversement.

\section{Existence, unicité et schéma d'Euler}

\begin{definition}[Équation différentielle stochastique, solutions forte et faible]\label{def:qm:stochastic-differential-equations:sde}
Une \emph{équation différentielle stochastique}\index{equation@équation différentielle stochastique} est
\[
dX_t = \mu(t, X_t)\,dt + \sigma(t, X_t)\,dW_t, \qquad X_0 = x_0,
\]
pour $\mu, \sigma$ mesurables. Une \emph{solution forte}\index{solution forte} sur un espace de probabilité donné
muni d'un \omterm{def:qm:brownian-motion:bm}{mouvement brownien} $W$ donné est un processus continu adapté vérifiant
$X_t = x_0 + \int_0^t\mu(s, X_s)\,ds + \int_0^t\sigma(s, X_s)\,dW_s$. Une \emph{solution
faible}\index{solution faible} est un couple $(X, W)$ sur un espace de probabilité filtré quelconque vérifiant
l'équation~: le \omterm{def:qm:brownian-motion:bm}{mouvement brownien} fait partie de la réponse.
\end{definition}

\begin{theorem}[Existence et unicité]\label{thm:qm:stochastic-differential-equations:exist}
Si $|\mu(t,x) - \mu(t,y)| + |\sigma(t,x) - \sigma(t,y)| \le K|x - y|$ et $|\mu(t,x)| + |\sigma(t,x)|
\le K(1 + |x|)$, l'équation admet une unique \omterm{def:qm:stochastic-differential-equations:sde}{solution forte}, avec $\E[\sup_{t\le T}X_t^2] <
\infty$.
\end{theorem}

\begin{proof}[Preuve partielle]
Itération de Picard $X^{(n+1)}_t = x_0 + \int_0^t\mu(X^{(n)})\,ds + \int_0^t\sigma(X^{(n)})\,dW$~: par
l'isométrie et l'inégalité de Doob, $\E[\sup_{s\le t}|X^{(n+1)}_s - X^{(n)}_s|^2] \le C\int_0^t
\E[\sup_{u\le s}|X^{(n)}_u - X^{(n-1)}_u|^2]\,ds$, donc les écarts sont bornés par $C^nt^n/n!$
et les itérées convergent~; le lemme de Gronwall donne l'unicité (Karatzas et Shreve, 1991, §5.2).
\end{proof}

Le coefficient en racine carrée $\eta\sqrt v$ n'est pas lipschitzien en zéro, ce qui explique que le
\omterm{def:qm:stochastic-differential-equations:sqrt}{processus en racine carrée} demande sa propre théorie ci-dessous. Les \omterm{def:qm:stochastic-differential-equations:sde}{solutions faibles} comptent, car les changements de mesure
(chapitre~5) les produisent~: l'équation $dX = \operatorname{sign}(X)\,dW$ admet une \omterm{def:qm:stochastic-differential-equations:sde}{solution faible}, un
\omterm{def:qm:brownian-motion:bm}{mouvement brownien}, mais aucune \omterm{def:qm:stochastic-differential-equations:sde}{solution forte}.

\begin{definition}[Schéma d'Euler--Maruyama]\label{def:qm:stochastic-differential-equations:euler}
Le \emph{schéma d'Euler--Maruyama}\index{schéma d'Euler--Maruyama} sur la grille $t_k = k\Delta t$ est
$\hat X_{k+1} = \hat X_k + \mu(t_k, \hat X_k)\Delta t + \sigma(t_k, \hat X_k)\sqrt{\Delta t}\,Z_k$, avec $Z_k$
indépendantes de loi normale centrée réduite.
\end{definition}

C'est l'\omterm{def:qm:ito-calculus:integral}{intégrale d'Itô} discrète du \cref{ch:qm:ito-calculus}~: les coefficients sont figés au
point de gauche de chaque intervalle. Sous les hypothèses lipschitziennes, elle converge lorsque $\Delta t \to 0$, trajectoire par trajectoire à l'ordre
$\tfrac12$ et en loi à l'ordre 1 (chapitre~26 précise les deux ordres). Sa faiblesse est
qu'elle ignore la géométrie de l'espace d'état~: un incrément gaussien peut faire passer un processus positif
sous zéro.

\section{Les diffusions de référence}

\begin{definition}[Mouvement brownien géométrique]\label{def:qm:stochastic-differential-equations:gbm}
Un \emph{mouvement brownien géométrique}\index{mouvement brownien géométrique} vérifie $dS = \mu S\,dt + \sigma
S\,dW$~; par l'\cref{ex:qm:ito-calculus:log}, $S_t = S_0\exp((\mu - \tfrac12\sigma^2)t + \sigma W_t)$, de loi lognormale.
\end{definition}

\begin{definition}[Retour à la moyenne, processus d'Ornstein--Uhlenbeck, demi-vie]\label{def:qm:stochastic-differential-equations:ou}
Un processus présente un \emph{retour à la moyenne}\index{retour à la moyenne} lorsque sa dérive le ramène vers un niveau.
Le \emph{processus d'Ornstein--Uhlenbeck}\index{processus d'Ornstein--Uhlenbeck} est le cas gaussien
$dX = \kappa(\bar x - X)\,dt + \sigma\,dW$, $\kappa > 0$. La \emph{demi-vie}\index{demi-vie} d'un
processus à retour à la moyenne est le temps $\ln 2/\kappa$ au bout duquel l'écart moyen au
niveau a été réduit de moitié.
\end{definition}

\begin{proposition}[La transition d'Ornstein--Uhlenbeck]\label{prop:qm:stochastic-differential-equations:ou}
$X_t = \bar x + (X_0 - \bar x)e^{-\kappa t} + \sigma\int_0^te^{-\kappa(t-s)}\,dW_s$~; sachant $X_0$, $X_t$ suit
une loi normale de moyenne $\bar x + (X_0 - \bar x)e^{-\kappa t}$ et de variance $\sigma^2(1 - e^{-2\kappa t})/(2\kappa)$,
qui tend vers $\mathcal N(\bar x, \sigma^2/2\kappa)$.
\end{proposition}

\begin{proof}
La formule d'Itô appliquée à $e^{\kappa t}(X_t - \bar x)$ donne $d(e^{\kappa t}(X_t - \bar x)) = \sigma e^{\kappa t}\,dW_t$~;
on intègre, puis on utilise l'isométrie pour la variance de l'intégrale de Wiener.
\end{proof}

La transition exacte permet de simuler le \omterm{def:qm:stochastic-differential-equations:ou}{processus d'Ornstein--Uhlenbeck} à n'importe quel pas de temps
sans erreur (la \cref{fig:qm:stochastic-differential-equations:paths}), et son échantillonnage discret
est une autorégression de coefficient $e^{-\kappa\Delta t}$ (chapitre~17). Les spreads, les bases de
financement et, dans le livre~6, les taux courts sont modélisés de cette façon.

\begin{definition}[Processus en racine carrée, condition de Feller]\label{def:qm:stochastic-differential-equations:sqrt}
Le \emph{processus en racine carrée}\index{processus en racine carrée} est $dv = \kappa(\bar v - v)\,dt + \eta\sqrt
v\,dW$ avec $\kappa, \bar v, \eta > 0$ et $v_0 > 0$. La \emph{condition
de Feller}\index{condition de Feller} est $2\kappa\bar v \ge \eta^2$~; le rapport $2\kappa\bar v/\eta^2$ s'appelle le ratio de Feller.
\end{definition}

\begin{theorem}[Feller]\label{thm:qm:stochastic-differential-equations:feller}
Le \omterm{def:qm:stochastic-differential-equations:sqrt}{processus en racine carrée} admet une unique \omterm{def:qm:stochastic-differential-equations:sde}{solution forte} à valeurs positives ou nulles. Il n'atteint jamais zéro si la
\omterm{def:qm:stochastic-differential-equations:sqrt}{condition de Feller} est vérifiée, et il atteint zéro avec une probabilité positive (et est instantanément réfléchi)
si elle n'est pas vérifiée. Sachant $v_s$, $v_t = c\,Y$ avec $Y$ suivant une loi du khi-deux non centrée à $d = 4\kappa\bar
v/\eta^2$ degrés de liberté et de paramètre de non-centralité $v_se^{-\kappa(t-s)}/c$, où $c = \eta^2(1 -
e^{-\kappa(t-s)})/(4\kappa)$~; sa \omterm{def:qm:stochastic-differential-equations:kolmogorov}{loi stationnaire} est $\mathrm{Gamma}(2\kappa\bar v/\eta^2,\ \eta^2/2\kappa)$.
\end{theorem}

\admitted

L'unicité repose sur l'argument de Yamada--Watanabe pour des coefficients höldériens d'exposant $\tfrac12$~; la
classification de la frontière est due à Feller (1951), la loi de transition figure chez Cox, Ingersoll et Ross (1985).
Le \omterm{def:qm:stochastic-differential-equations:sqrt}{processus en racine carrée} est la variance du modèle de Heston et le taux court du modèle CIR
(les livres~5 et~6)~; les ajustements à volatilité stochastique sur les sourires de volatilité actions demandent en général
une forte volatilité de la volatilité $\eta$ et violent donc la condition. Le processus du desk a $\kappa = 2$, $\bar v =
v_0 = 0.04$ et $\eta = 0.6$~: un ratio de Feller égal à 0.44, une \omterm{def:qm:stochastic-differential-equations:kolmogorov}{loi stationnaire} de moyenne 0.04 et
écart-type 0.06, et 27.9~\% de sa masse stationnaire est inférieure à un dixième de la moyenne.

\begin{omfigure}
\begin{tikzpicture}
\begin{groupplot}[group style={group size=2 by 1,horizontal sep=1.7cm},
  width=7.6cm,height=4.8cm,xlabel={$t$ (années)},xmin=0,xmax=1,
  tick label style={font=\small},label style={font=\small},grid=major,grid style={gray!25},
  table/col sep=comma]
\nextgroupplot[ylabel={$X_t$},ymin=-0.35,ymax=0.6,
  yticklabel style={/pgf/number format/fixed}]
\addplot[omDef,thin] table[x=t,y=ou0]{figdata/methods/04-stochastic-differential-equations/paths.csv};
\addplot[omThm,thin] table[x=t,y=ou1]{figdata/methods/04-stochastic-differential-equations/paths.csv};
\addplot[omProp,thin] table[x=t,y=ou2]{figdata/methods/04-stochastic-differential-equations/paths.csv};
\addplot[black,thick,dashed] table[x=t,y=ou_mean]{figdata/methods/04-stochastic-differential-equations/paths.csv};
\nextgroupplot[ylabel={$v_t$},ymin=0,ymax=0.11,
  yticklabel style={/pgf/number format/fixed}]
\addplot[omDef,thin] table[x=t,y=v0]{figdata/methods/04-stochastic-differential-equations/paths.csv};
\addplot[omThm,thin] table[x=t,y=v1]{figdata/methods/04-stochastic-differential-equations/paths.csv};
\addplot[omProp,thin] table[x=t,y=v2]{figdata/methods/04-stochastic-differential-equations/paths.csv};
\addplot[black,dashed,domain=0:1] {0.04};
\end{groupplot}
\end{tikzpicture}

\omcaption{À gauche~: trois trajectoires d'Ornstein--Uhlenbeck partant de 0.5 avec $\kappa = 2$ (\omterm{def:qm:stochastic-differential-equations:ou}{demi-vie} 0.35 an),
$\bar x = 0$, $\sigma = 0.2$, et leur trajectoire moyenne $0.5e^{-2t}$ (tirets). À droite~: trois trajectoires du
\omterm{def:qm:stochastic-differential-equations:sqrt}{processus en racine carrée} du desk ($\kappa = 2$, $\bar v = 0.04$, $\eta = 0.6$), échantillonnées exactement chaque
jour~; avec un ratio de Feller égal à 0.44, elles passent de longues périodes près de zéro. Données~: tutoriel du
chapitre, graine fixée.}\label{fig:qm:stochastic-differential-equations:paths}
\end{omfigure}

Le schéma d'Euler élémentaire à partir d'une petite valeur de $v$ est gaussien, de moyenne $v + \kappa(\bar v - v)\Delta t$ et
d'écart-type $\eta\sqrt{v\Delta t}$~: à partir de $v = 0.004$, un pas quotidien devient négatif avec
une probabilité de 3.6~\%, et une trajectoire passe de nombreux jours près de zéro.
La \Cref{fig:qm:stochastic-differential-equations:feller} compte les trajectoires qui produisent une
variance négative au cours d'une année~: 75~\% aux paramètres du desk, et, ce qui surprend le plus, les
mêmes 75~\% avec quatre ou seize pas par jour. Raffiner ne peut pas aider, car le processus exact
lui-même s'approche arbitrairement de zéro sur ces trajectoires~; seul un schéma qui respecte la
frontière le peut. Même avec un ratio de Feller exactement égal à un, un quart des trajectoires d'Euler quotidiennes
devient négatif~: la condition protège le processus, non le schéma.

\begin{omfigure}
\begin{tikzpicture}
\begin{axis}[width=11cm,height=5cm,xmode=log,log basis x=2,
  xlabel={ratio de Feller $2\kappa\bar v/\eta^2$ ($\kappa = 2$, $\bar v = 0.04$)},ylabel={trajectoires avec un $v$ négatif},
  tick label style={font=\small},label style={font=\small},
  xmin=0.14,xmax=4.5,ymin=0,ymax=1,grid=major,grid style={gray!25},
  xtick={0.25,0.5,1,2,4},xticklabels={0.25,0.5,1,2,4},
  legend columns=2,legend style={font=\footnotesize,at={(0.5,-0.3)},anchor=north,draw=none,fill=none,/tikz/every even column/.append style={column sep=8pt}},
  table/col sep=comma]
\addplot[omThm,very thick,mark=*,mark size=1.6pt] table[x=ratio,y=daily]{figdata/methods/04-stochastic-differential-equations/feller.csv};
\addplot[omDef,thick,dashed,mark=triangle*,mark size=1.8pt] table[x=ratio,y=fine]{figdata/methods/04-stochastic-differential-equations/feller.csv};
\addplot[gray,thin] coordinates {(1,0) (1,1)};
\addplot[black,only marks,mark=o,mark size=3pt] coordinates {(0.4444,0.752)};
\node[font=\footnotesize,anchor=west] at (axis cs:0.5,0.93) {desk~: 0.44, 75~\%};
\legend{pas quotidiens,quatre pas par jour}
\end{axis}
\end{tikzpicture}

\omcaption{Part des trajectoires du schéma d'Euler élémentaire du \omterm{def:qm:stochastic-differential-equations:sqrt}{processus en racine carrée} qui produisent une
variance négative au cours d'une année, en fonction du ratio de Feller (volatilité de la volatilité $\eta$ de 0.2 à 1.0). Raffiner le
pas ne supprime pas les échecs~; les schémas exact et de troncature complète n'en ont aucun. Données~: 20\,000
et 10\,000 trajectoires simulées avec graine fixée par point.}\label{fig:qm:stochastic-differential-equations:feller}
\end{omfigure}

Deux correctifs sont classiques. Le schéma exact échantillonne la transition du chi-deux non central
énoncée au \cref{thm:qm:stochastic-differential-equations:feller} et reste positif par construction.
\emph{Troncature complète} (Lord, Koekkoek et van Dijk, 2010) conserve le pas d'Euler mais utilise $v^+ =
\max(v, 0)$ dans les deux coefficients, en rapportant $v^+$~: elle ne prend jamais la racine carrée d'un
nombre négatif, et son biais s'annule lorsque $\Delta t \to 0$. Le livre~5, chapitre~23, ajoute le
schéma quadratique-exponentiel utilisé pour Heston en production.

\section{Générateurs et équations de Kolmogorov}

\begin{definition}[Générateur infinitésimal]\label{def:qm:stochastic-differential-equations:generator}
Le \emph{générateur infinitésimal}\index{générateur infinitésimal} d'une
diffusion homogène en temps $dX = \mu(X)\,dt + \sigma(X)\,dW$ est l'opérateur $\mathcal Lf(x) = \lim_{t\downarrow 0}
(\E_x[f(X_t)] - f(x))/t = \mu(x)f'(x) + \tfrac12\sigma^2(x)f^{\prime\prime}(x)$ sur les fonctions de classe $C^2$.
\end{definition}

\begin{proposition}[Formule de Dynkin]\label{prop:qm:stochastic-differential-equations:dynkin}
Pour $f \in C^2$ à support compact et un \omterm{def:qm:probability-at-speed:stopping}{temps d'arrêt} $\tau$ avec $\E_x[\tau] < \infty$,
$\E_x[f(X_\tau)] = f(x) + \E_x\int_0^\tau\mathcal Lf(X_s)\,ds$.
\end{proposition}

\begin{proof}
D'après la formule d'Itô, $f(X_t) - f(x) - \int_0^t\mathcal Lf(X_s)\,ds = \int_0^tf'(X_s)\sigma(X_s)\,dW_s$, une
\omterm{def:qm:probability-at-speed:martingale}{martingale} à intégrande borné~; appliquer le théorème d'arrêt optionnel en $\tau \wedge n$ et faire tendre $n \to \infty$.
\end{proof}

\begin{definition}[Équations de Kolmogorov, loi stationnaire]\label{def:qm:stochastic-differential-equations:kolmogorov}
Pour $u(t, x) = \E[g(X_T) \mid X_t = x]$, l'\emph{équation de Kolmogorov rétrograde}\index{équation de Kolmogorov rétrograde}
est $\partial_tu + \mathcal Lu = 0$ avec $u(T, \cdot) = g$. La densité de transition
$p(t, y)$ de $X_t$ résout l'\emph{équation de Kolmogorov progressive}\index{équation de Kolmogorov progressive},
aussi appelée \emph{équation de Fokker--Planck}\index{équation de Fokker--Planck}, $\partial_tp =
\mathcal L^*p = -\partial_y(\mu p) + \tfrac12\partial_{yy}(\sigma^2p)$. Une \emph{loi stationnaire}\index{loi stationnaire} d'un \omterm{def:qm:brownian-motion:markov}{processus de Markov}
est une loi qui, prise comme loi de
$X_0$, est la loi de chaque $X_t$~; pour une diffusion, sa densité résout $\mathcal L^*p = 0$.
\end{definition}

L'équation rétrograde regarde un paiement depuis la position où l'on se trouve~; l'équation progressive fait avancer
une densité dans le temps. Pour le \omterm{def:qm:stochastic-differential-equations:ou}{processus d'Ornstein--Uhlenbeck}, $\mathcal L^*p = 0$ s'intègre une fois en
$\kappa(\bar x - y)p = \tfrac12\sigma^2p'$, dont la solution est la densité de $\mathcal N(\bar x, \sigma^2/2\kappa)$~; pour le
\omterm{def:qm:stochastic-differential-equations:sqrt}{processus en racine carrée}, $\kappa(\bar v - y)p = \tfrac12\eta^2(yp)'$ donne $p \propto y^{2\kappa\bar
v/\eta^2 - 1}e^{-2\kappa y/\eta^2}$, la loi Gamma du \cref{thm:qm:stochastic-differential-equations:feller}.
Lorsque le ratio de Feller est inférieur à un, l'exposant est négatif et la densité est infinie
en zéro (\cref{fig:qm:stochastic-differential-equations:stationary}).

\begin{omfigure}
\begin{tikzpicture}
\begin{groupplot}[group style={group size=2 by 1,horizontal sep=1.7cm},
  width=7.6cm,height=4.8cm,xlabel={$v$},xmin=0,xmax=0.16,ymin=0,
  tick label style={font=\small},label style={font=\small},grid=major,grid style={gray!25},
  xticklabel style={/pgf/number format/fixed},scaled x ticks=false,
  table/col sep=comma]
\nextgroupplot[ylabel={densité},ymax=66,title={\small desk~: ratio de Feller 0.44}]
\addplot[omDef,thick,const plot mark mid] table[x=v,y=desk_hist]{figdata/methods/04-stochastic-differential-equations/stationary.csv};
\addplot[omThm,only marks,mark=*,mark size=1.3pt] table[x=v,y=desk_pdf]{figdata/methods/04-stochastic-differential-equations/stationary.csv};
\nextgroupplot[ymax=26,title={\small $\eta = 0.2$~: ratio de Feller 4}]
\addplot[omDef,thick,const plot mark mid] table[x=v,y=f2_hist]{figdata/methods/04-stochastic-differential-equations/stationary.csv};
\addplot[omThm,only marks,mark=*,mark size=1.3pt] table[x=v,y=f2_pdf]{figdata/methods/04-stochastic-differential-equations/stationary.csv};
\end{groupplot}
\end{tikzpicture}

\omcaption{\omterm{def:qm:stochastic-differential-equations:kolmogorov}{Lois stationnaires} du \omterm{def:qm:stochastic-differential-equations:sqrt}{processus en racine carrée} avec $\kappa = 2$, $\bar v = 0.04$~: histogrammes de
50\,000 trajectoires exactes après cinq ans (pas) comparées à la loi Gamma qui résout $\mathcal L^*p = 0$,
moyennée sur chaque classe (points). Avec le $\eta = 0.6$ du desk, la densité s'accumule en zéro~; avec
$\eta = 0.2$, zéro est inatteignable. Données~: le tutoriel du chapitre, avec graine fixée.}\label{fig:qm:stochastic-differential-equations:stationary}
\end{omfigure}

\section{Feynman--Kac}

\begin{theorem}[Feynman--Kac]\label{thm:qm:stochastic-differential-equations:fk}
Soit $X$ solution de l'équation de la \cref{def:qm:stochastic-differential-equations:sde}, $r \ge 0$ et $g$
continue et à croissance polynomiale, et soit $u \in C^{1,2}$ à croissance polynomiale solution de
\[
\partial_tu + \mu\,\partial_xu + \tfrac12\sigma^2\partial_{xx}u - r(t,x)u = 0, \qquad u(T, x) = g(x).
\]
Alors $u(t, x) = \E\bigl[e^{-\int_t^Tr(s, X_s)\,ds}g(X_T) \bigm| X_t = x\bigr]$.
\end{theorem}

\begin{proof}
Avec $D_s = e^{-\int_t^sr(u, X_u)\,du}$, la règle du produit d'Itô donne $d(D_su(s, X_s)) = D_s(\partial_su +
\mathcal Lu - ru)\,ds + D_s\sigma\partial_xu\,dW_s = D_s\sigma\partial_xu\,dW_s$. Les conditions de croissance rendent
l'intégrale stochastique une \omterm{def:qm:probability-at-speed:martingale}{martingale}, donc $u(t, x) = \E[D_Tu(T, X_T) \mid X_t = x]$.
\end{proof}

Le théorème vaut dans les deux sens~: une espérance se calcule en résolvant une EDP (chapitre~27), et une
EDP se résout en simulant des trajectoires (chapitre~26). La \cref{fig:qm:stochastic-differential-equations:map} rassemble les liens.
Avec $X = r$ lui-même un \omterm{def:qm:stochastic-differential-equations:ou}{processus d'Ornstein--Uhlenbeck} pour le taux court, et $g = 1$, $u$ est le
prix d'une obligation zéro-coupon, et l'EDP admet la solution affine $u = e^{A(\tau) - B(\tau)r}$ avec
$B(\tau) = (1 - e^{-\kappa\tau})/\kappa$ et $A(\tau) = (\bar r - \sigma^2/2\kappa^2)(B - \tau) - \sigma^2B^2/4\kappa$,
obtenue en substituant et en identifiant les termes en $r$. Avec $r_0 = 3\%$, $\bar r = 4\%$, $\kappa = 0.5$
et $\sigma = 1\%$, le prix à cinq ans vaut 0.8343, et 20\,000 trajectoires simulées donnent $0.8341 \pm 0.0002$.
Le livre~6, chapitre~7, transforme ce calcul en un modèle de la courbe.

\begin{omfigure}
\begin{tikzpicture}[font=\small,node distance=0.7cm and 2.7cm,
  bx/.style={draw=omDef,rounded corners=2pt,fill=omDef!6,inner sep=4pt,align=center,text width=4.2cm}]
\node[bx] (sde) {EDS\\$dX = \mu\,dt + \sigma\,dW$};
\node[bx,right=of sde] (gen) {générateur\\$\mathcal Lf = \mu f' + \tfrac12\sigma^2f^{\prime\prime}$};
\node[bx,below=of gen] (bwd) {équation rétrograde\\$\partial_tu + \mathcal Lu - ru = 0$, $u(T) = g$};
\node[bx,below=of sde] (exp) {espérance\\$u = \E_{t,x}[e^{-\int r}g(X_T)]$};
\node[bx,below=of bwd] (fwd) {équation progressive\\$\partial_tp = \mathcal L^*p$\\stationnaire~: $\mathcal L^*p = 0$};
\node[bx,below=of exp] (law) {loi de $X_t$\\densité $p(t, y)$};
\draw[-{Stealth},thick] (sde) -- node[above,font=\footnotesize] {Itô} (gen);
\draw[-{Stealth},thick] (gen) -- (bwd);
\draw[{Stealth}-{Stealth},thick,omThm] (bwd) -- node[above,font=\footnotesize,omThm] {Feynman--Kac} (exp);
\draw[-{Stealth},thick] (gen.east) to[bend left=35] node[right,font=\footnotesize] {adjoint} (fwd.east);
\draw[{Stealth}-{Stealth},thick] (fwd) -- (law);
\draw[-{Stealth},thick,gray] (sde) -- node[left,font=\footnotesize] {simuler} (exp);
\end{tikzpicture}

\omcaption{Comment les équations d'une diffusion s'articulent entre elles. Le générateur découle de la formule d'Itô~;
les espérances de gains résolvent l'équation rétrograde (Feynman--Kac, avec actualisation), les densités
résolvent l'équation progressive, et les \omterm{def:qm:stochastic-differential-equations:kolmogorov}{lois stationnaires} résolvent $\mathcal L^*p = 0$.}\label{fig:qm:stochastic-differential-equations:map}
\end{omfigure}

\section{Tutoriel~: le NaN de la nuit}\label{tut:qm:stochastic-differential-equations:nan}

\begin{tutorial}
\textbf{Objectif.} Reproduire la panne de la nuit, la mesurer à l'aune du ratio de Feller, la corriger de deux
façons, et vérifier la \omterm{def:qm:stochastic-differential-equations:kolmogorov}{loi stationnaire} et un prix de Feynman--Kac. \textbf{État final~:}
\cref{fig:qm:stochastic-differential-equations:paths,fig:qm:stochastic-differential-equations:feller,fig:qm:stochastic-differential-equations:stationary}
et les chiffres du problème du week-end.
\begin{tutsteps}
\item \textbf{Les schémas.} Le projet en cours échantillonne le \omterm{def:qm:stochastic-differential-equations:sqrt}{processus en racine carrée} exactement et par
Euler simple ou avec troncature complète.
\omcode{code/firm/mcengine/firm_mcengine.py}{150}{180}{Pas exacts et pas d'Euler du processus en racine carrée.}\label{lst:qm:stochastic-differential-equations:sqrt}
\item \textbf{Le décompte}~: la part des trajectoires d'Euler simple qui produisent une variance négative, et une
vérification de Feynman--Kac par simulation.
\omcode{code/methods/04-stochastic-differential-equations/python/qm_sde.py}{77}{90}{Prix d'une obligation zéro-coupon sous un taux court d'Ornstein--Uhlenbeck, par simulation et par l'EDP.}\label{lst:qm:stochastic-differential-equations:fk}
\item \textbf{Exécuter} \texttt{problem()}, \texttt{feller\_table()}, \texttt{stationary\_histograms()} et
\texttt{fig\_sde.py}~; les versions C++20 et Rust des pas d'Euler et d'Ornstein--Uhlenbeck se trouvent dans
\texttt{code/firm/mcengine/}.
\end{tutsteps}
\textbf{Pour aller plus loin.} Ajouter le schéma par réflexion ($v \leftarrow |v|$) et comparer sa moyenne sur un an
avec celle du schéma exact~; donner au processus de variance un niveau $\bar v(t)$ dépendant du temps et
vérifier la prédiction du générateur pour $\E[v_t]$.
\end{tutorial}

\section{Construction~: le moteur de Monte-Carlo, étape deux}\label{bld:qm:stochastic-differential-equations:mcengine}

\begin{build}
\bfield{Rôle} Faire progresser les diffusions que simule chaque chapitre suivant, avec le schéma dont chacune
a besoin, et ne jamais renvoyer de NaN.
\bfield{Interface} \texttt{euler\_maruyama(mu, sigma, x0, T, n\_steps, n\_paths, seed)}~;
\texttt{ou\_exact(x0, kappa, xbar, sigma, T, n\_steps, n\_paths, seed)}~; \texttt{sqrt\_exact(v0,
kappa, vbar, eta, \dots)}~; \texttt{sqrt\_euler(\dots, scheme)} avec \texttt{plain} ou
\texttt{full\_truncation}~; \texttt{feller\_ratio(kappa, vbar, eta)}. C++20 et Rust~:
\texttt{ou\_exact\_path} et \texttt{sqrt\_euler\_path} sur le \texttt{NormalStream} de l'étape un.
\bfield{Règles} Paramètres dans les notations de la série ($\kappa$, $\bar x$, $\bar v$, $\eta$)~; transitions exactes là où
elles existent~; un schéma susceptible de sortir de l'espace d'états est désigné comme tel et n'est jamais le choix par défaut.
\bfield{Tests d'acceptation} Moyenne et variance stationnaires des deux processus~; Euler converge vers la
moyenne exacte du \omterm{def:qm:stochastic-differential-equations:ou}{processus d'Ornstein--Uhlenbeck}~; le schéma exact en racine carrée est positif ou nul, avec
les bons moments~; Euler simple échoue aux paramètres du desk et la troncature complète n'échoue jamais, dans les trois
langages.
\bfield{Pour aller plus loin} Le schéma quadratique-exponentiel~; un pas de Milstein (chapitre~26).
\end{build}

\begin{omsources}
\item W.~Feller, ``Two singular diffusion problems'', \emph{Annals of Mathematics} 54, 1951.
\item J.~C.~Cox, J.~E.~Ingersoll and S.~A.~Ross, ``A theory of the term structure of interest
rates'', \emph{Econometrica} 53, 1985.
\item G.~E.~Uhlenbeck and L.~S.~Ornstein, ``On the theory of the Brownian motion'',
\emph{Physical Review} 36, 1930.
\item M.~Kac, ``On distributions of certain Wiener functionals'', \emph{Transactions of the AMS}
65, 1949.
\item R.~Lord, R.~Koekkoek and D.~van Dijk, ``A comparison of biased simulation schemes for
stochastic volatility models'', \emph{Quantitative Finance} 10, 2010.
\end{omsources}

\section{Exercices}

\begin{exercise}[$\star$]\label{exo:qm:stochastic-differential-equations:1}
Un spread suit un \omterm{def:qm:stochastic-differential-equations:ou}{processus d'Ornstein--Uhlenbeck} de paramètre $\kappa = 2$ par an. Quelle est sa \omterm{def:qm:stochastic-differential-equations:ou}{demi-vie}
en années et en jours de négociation~?
\end{exercise}

\begin{exercise}[$\star$]\label{exo:qm:stochastic-differential-equations:2}
Une action suit un \omterm{def:qm:stochastic-differential-equations:gbm}{mouvement brownien géométrique} de paramètres $\mu = 5\%$ et $\sigma = 20\%$. Quelle est la probabilité
qu'elle soit inférieure à son prix initial au bout d'un an~?
\end{exercise}

\begin{exercise}[$\star$]\label{exo:qm:stochastic-differential-equations:3}
Écrire le générateur du \omterm{def:qm:stochastic-differential-equations:ou}{processus d'Ornstein--Uhlenbeck}, l'appliquer à $f(x) = x$, et en déduire l'EDO
satisfaite par $\E[X_t]$.
\end{exercise}

\begin{exercise}[$\star\star$]\label{exo:qm:stochastic-differential-equations:4}
À partir de l'\omterm{def:qm:stochastic-differential-equations:kolmogorov}{équation de Kolmogorov progressive}, déterminer la variance stationnaire du \omterm{def:qm:stochastic-differential-equations:ou}{processus
d'Ornstein--Uhlenbeck} avec $\kappa = 2$, $\sigma = 0.4$.
\end{exercise}

\begin{exercise}[$\star\star$]\label{exo:qm:stochastic-differential-equations:5}
À l'aide de la formule de Dynkin, calculer $\E[v_{0.5}]$ pour le \omterm{def:qm:stochastic-differential-equations:sqrt}{processus en racine carrée} avec $v_0 = 0.09$, $\bar v =
0.04$, $\kappa = 2$.
\end{exercise}

\begin{exercise}[$\star\star$]\label{exo:qm:stochastic-differential-equations:6}
Vérifier le prix zéro-coupon à cinq ans du chapitre (0.8343) à partir de la formule de $A$ et $B$.
\end{exercise}

\begin{exercise}[$\star\star\star$]\label{exo:qm:stochastic-differential-equations:7}
\emph{Programmation.} Avec \texttt{negative\_fraction}, déterminer la part des trajectoires d'Euler simples quotidiennes qui
deviennent négatives au cours d'une année lorsque le rapport de Feller vaut exactement un ($\eta = 0.4$), ainsi que le plus grand $\eta$ du
tableau pour lequel elle est inférieure à 1~\%.
\end{exercise}

\begin{exercise}[$\star\star\star$]\label{exo:qm:stochastic-differential-equations:8}
\emph{Trouver l'erreur.} «~On tronque la variance à zéro après chaque pas d'Euler, donc le schéma est
désormais exact.~»
\end{exercise}

\section{Problème~: le NaN de la nuit}

\begin{problem}[{Problème du week-end --- un processus de variance qui viole la condition de Feller}]\label{pb:qm:stochastic-differential-equations:1}
Le processus de variance du desk est $dv = \kappa(\bar v - v)\,dt + \eta\sqrt v\,dW$ avec $v_0 = \bar v = 0.04$,
$\kappa = 2$ et $\eta = 0.6$, simulé sur une année de 252 jours.

\medskip\noindent\textbf{Partie I --- Le processus.}
\begin{enumerate}
\item Quel est le rapport de Feller~? Le zéro est-il atteignable~?
\item Quelles sont la \omterm{def:qm:stochastic-differential-equations:kolmogorov}{loi stationnaire}, sa moyenne et son écart-type~?
\item Quelle part de la masse stationnaire se situe en dessous de 0.004~?
\item Quelle volatilité de la volatilité rendrait le rapport de Feller exactement égal à un~?
\item Quelle est la \omterm{def:qm:stochastic-differential-equations:ou}{demi-vie} du retour de la variance moyenne vers $\bar v$~?
\end{enumerate}

\medskip\noindent\textbf{Partie II --- Euler simple.}
\begin{enumerate}[resume]
\item À partir de $v = 0.004$, quelle est la probabilité qu'un pas d'Euler quotidien devienne négatif~?
\item Quelle part des trajectoires produit une variance négative au cours de l'année~?
\item Avec quatre et seize pas par jour, et avec des pas hebdomadaires~?
\item Pourquoi raffiner le pas n'aide-t-il pas~?
\item Que fait un seul NaN à un calcul de risque qui somme sur les trajectoires~?
\end{enumerate}

\medskip\noindent\textbf{Partie III --- Corrections.}
\begin{enumerate}[resume]
\item Quels sont le minimum et la moyenne annuelle donnés par le schéma exact~?
\item Quelle moyenne annuelle donne la troncature complète, et sur quelle part des pas vaut-elle zéro~?
\item Lequel des deux est biaisé, et comment le biais se comporte-t-il lorsque $\Delta t \to 0$~?
\item Que ferait la réflexion ($v \leftarrow |v|$) sur la moyenne au voisinage de zéro~?
\item Pourquoi les moteurs Heston de production utilisent-ils un autre schéma~?
\end{enumerate}

\medskip\noindent\textbf{Partie IV --- Jugement.}
\begin{enumerate}[resume]
\item Le desk doit-il imposer la \omterm{def:qm:stochastic-differential-equations:sqrt}{condition de Feller} dans sa calibration~?
\item Que doit vérifier un test automatisé du moteur de simulation~?
\item Que faut-il consigner dans la note de revue de modèle portant sur cet incident~?
\item Énoncer le \emph{résultat nommé}~: le rapport de Feller et la part des trajectoires défaillantes.
\item En une phrase~: que protège la \omterm{def:qm:stochastic-differential-equations:sqrt}{condition de Feller}, et que ne protège-t-elle pas~?
\end{enumerate}
\end{problem}

\section{Questions d'entretien}

\begin{interviewq}[$\star$ \iqroles{trader, chercheur}]\label{iq:qm:stochastic-differential-equations:1}
Résolvez $dS = \mu S\,dt + \sigma S\,dW$. Quelle est la médiane de $S_T$~?
\end{interviewq}

\begin{interviewq}[$\star$ \iqroles{chercheur}]\label{iq:qm:stochastic-differential-equations:2}
Un spread revient vers sa moyenne avec une \omterm{def:qm:stochastic-differential-equations:ou}{demi-vie} de 10 jours. Quelle est la valeur de $\kappa$, et quelle part de
l'écart actuel devrait subsister après 30 jours~?
\end{interviewq}

\begin{interviewq}[$\star\star$ \iqroles{chercheur, banque}]\label{iq:qm:stochastic-differential-equations:3}
Qu'est-ce que la \omterm{def:qm:stochastic-differential-equations:sqrt}{condition de Feller}, et que se passe-t-il pour un \omterm{def:qm:stochastic-differential-equations:sqrt}{processus en racine carrée} et pour sa simulation
lorsqu'elle n'est pas vérifiée~?
\end{interviewq}

\begin{interviewq}[$\star\star$ \iqroles{chercheur, banque}]\label{iq:qm:stochastic-differential-equations:4}
Énoncez la formule de Feynman--Kac et utilisez-la pour écrire l'EDP satisfaite par $\E[e^{-rT}g(S_T)]$ sous
un \omterm{def:qm:stochastic-differential-equations:gbm}{mouvement brownien géométrique}.
\end{interviewq}

\begin{interviewq}[$\star\star$ \iqroles{développeur}]\label{iq:qm:stochastic-differential-equations:5}
La simulation Monte Carlo de nuit renvoie parfois NaN. Comment en trouvez-vous la cause et
comment faites-vous pour qu'il ne puisse plus se produire~?
\end{interviewq}

\begin{interviewq}[$\star\star\star$ \iqroles{chercheur}]\label{iq:qm:stochastic-differential-equations:6}
Quelle est la différence entre une \omterm{def:qm:stochastic-differential-equations:sde}{solution forte} et une \omterm{def:qm:stochastic-differential-equations:sde}{solution faible}~? Donnez une équation qui admet l'une
mais pas l'autre.
\end{interviewq}
